\documentclass[10pt]{article}
\usepackage[UTF8]{ctex}
\usepackage{amsmath, amssymb}
\usepackage{geometry}
\geometry{a4paper, margin=1.5cm}
\usepackage{enumitem}

\title{导数入门——综合练习}
\author{}
\date{}

\begin{document}
\maketitle

\section*{致亲爱的同学}
欢迎进入导数的世界！这一份练习旨在帮助你从基础到灵活运用，一步步掌握导数的计算。每一道题都配有详细的思路点拨和解答步骤，请先独立思考，再对照参考答案，相信你会收获满满。

\section{基础题——直接套用公式和法则}
这些题目考查幂函数、指数函数、对数函数、三角函数的求导公式，以及和差法则。请先熟悉基本公式，再动手计算。

\begin{enumerate}[label=\arabic*.]
    \item 求 $f(x) = 5x^3 - 2x^2 + 3x - 7$ 的导数。
    
    \textbf{解：} 这是一个多项式函数，直接使用幂函数公式 $(x^n)' = n x^{n-1}$ 和和差法则逐项求导，常数项导数为 $0$。
    \begin{align*}
        f'(x) &= (5x^3)' - (2x^2)' + (3x)' - (7)' \\
        &= 5 \cdot 3x^{2} - 2 \cdot 2x^{1} + 3 \cdot 1x^{0} - 0 \\
        &= 15x^2 - 4x + 3.
    \end{align*}
    \textbf{答案：} $f'(x) = 15x^2 - 4x + 3$。

    \item 求 $f(x) = e^{2x} + \ln x$ 的导数。
    
    \textbf{解：} 两项相加，分别求导。第一项 $e^{2x}$ 是复合函数（外层 $e^u$，内层 $u=2x$），需用链式法则；第二项 $\ln x$ 直接用公式。
    \begin{align*}
        f'(x) &= (e^{2x})' + (\ln x)' \\
        &= e^{2x} \cdot (2x)' + \frac{1}{x} \\
        &= e^{2x} \cdot 2 + \frac{1}{x} \\
        &= 2e^{2x} + \frac{1}{x}.
    \end{align*}
    \textbf{答案：} $f'(x) = 2e^{2x} + \frac{1}{x}$。

    \item 求 $f(x) = 3\sin x - 4\cos x$ 的导数。
    
    \textbf{解：} 两项相减，分别用 $(\sin x)' = \cos x$ 和 $(\cos x)' = -\sin x$，常数因子可提出。
    \begin{align*}
        f'(x) &= 3(\sin x)' - 4(\cos x)' \\
        &= 3\cos x - 4(-\sin x) \\
        &= 3\cos x + 4\sin x.
    \end{align*}
    \textbf{答案：} $f'(x) = 3\cos x + 4\sin x$。
\end{enumerate}

\section{进阶题——需要组合法则}
当函数形式变为乘积、商或复合结构时，就需要灵活运用乘积法则、商法则和链式法则了。试试看！

\begin{enumerate}[label=\arabic*., resume]
    \item 求 $f(x) = x^2 e^x$ 的导数。
    
    \textbf{解：} 这是两个函数相乘，符合乘积法则的形式。设 $u=x^2$，$v=e^x$，则 $f' = u'v + uv'$。
    \begin{align*}
        f'(x) &= (x^2)' \cdot e^x + x^2 \cdot (e^x)' \\
        &= 2x \cdot e^x + x^2 \cdot e^x \\
        &= e^x (2x + x^2).
    \end{align*}
    \textbf{答案：} $f'(x) = e^x (x^2 + 2x)$。

    \item 求 $f(x) = \dfrac{x}{x^2+1}$ 的导数。
    
    \textbf{解：} 这是分式函数，符合商法则的形式。设 $u = x$，$v = x^2+1$，则 $f' = \dfrac{u'v - uv'}{v^2}$。
    \begin{align*}
        f'(x) &= \frac{(x)' (x^2+1) - x (x^2+1)'}{(x^2+1)^2} \\
        &= \frac{1 \cdot (x^2+1) - x \cdot (2x)}{(x^2+1)^2} \\
        &= \frac{x^2+1 - 2x^2}{(x^2+1)^2} \\
        &= \frac{1 - x^2}{(x^2+1)^2}.
    \end{align*}
    \textbf{答案：} $f'(x) = \dfrac{1 - x^2}{(x^2+1)^2}$。

    \item 求 $f(x) = \sin(2x) \cos(2x)$ 的导数。
    
    \textbf{解：} 有两种方法。这里我们用三角恒等式化简：$\sin(2x)\cos(2x) = \dfrac{1}{2}\sin(4x)$，然后求导。
    $$f(x) = \frac{1}{2}\sin(4x).$$
    令 $u=4x$，则 $f(x) = \frac{1}{2}\sin u$，由链式法则：
    $$f'(x) = \frac{1}{2} \cdot \cos(4x) \cdot (4x)' = \frac{1}{2} \cdot \cos(4x) \cdot 4 = 2\cos(4x).$$
    \textbf{答案：} $f'(x) = 2\cos(4x)$。
    
    \textit{补充：} 若直接用乘积法则，可得相同结果：$f'(x) = 2\cos^2(2x) - 2\sin^2(2x) = 2\cos(4x)$。

    \item 求 $f(x) = \ln(\sin x)$ 的导数。
    
    \textbf{解：} 这是复合函数，外层是 $\ln u$，内层是 $u = \sin x$，用链式法则。注意定义域要求 $\sin x > 0$。
    \begin{align*}
        f'(x) &= \frac{1}{\sin x} \cdot (\sin x)' \\
        &= \frac{1}{\sin x} \cdot \cos x = \cot x.
    \end{align*}
    \textbf{答案：} $f'(x) = \cot x$（在 $\sin x > 0$ 的区间内成立）。
\end{enumerate}

\section{挑战题——灵活运用}
综合以上法则，解决更复杂的函数求导问题。试试你的身手！

\begin{enumerate}[label=\arabic*., resume]
    \item 求 $f(x) = (x^3+1)^5 \cdot e^{2x}$ 的导数。
    
    \textbf{解：} 这是两个函数相乘：$u = (x^3+1)^5$，$v = e^{2x}$。需用乘积法则，而 $u$ 本身又是复合函数（幂函数嵌套多项式），需用链式法则。
    \begin{align*}
        u' &= 5(x^3+1)^4 \cdot (x^3+1)' = 5(x^3+1)^4 \cdot 3x^2 = 15x^2 (x^3+1)^4, \\
        v' &= e^{2x} \cdot 2 = 2e^{2x}.
    \end{align*}
    由乘积法则：
    \begin{align*}
        f'(x) &= u' v + u v' \\
        &= 15x^2 (x^3+1)^4 \cdot e^{2x} + (x^3+1)^5 \cdot 2e^{2x} \\
        &= e^{2x} (x^3+1)^4 \bigl[ 15x^2 + 2(x^3+1) \bigr] \\
        &= e^{2x} (x^3+1)^4 (15x^2 + 2x^3 + 2) \\
        &= e^{2x} (x^3+1)^4 (2x^3 + 15x^2 + 2).
    \end{align*}
    \textbf{答案：} $f'(x) = e^{2x} (x^3+1)^4 (2x^3 + 15x^2 + 2)$。
    
    \textit{技巧：} 遇到多个因子相乘时，可以先分别求导，再代入乘积法则，最后提取公因式化简。

    \item 求 $f(x) = \sqrt{x} + \dfrac{1}{\sqrt{x}}$ 的导数。
    
    \textbf{解：} 先将根式写成幂的形式：$\sqrt{x} = x^{1/2}$，$\dfrac{1}{\sqrt{x}} = x^{-1/2}$，然后使用幂函数公式。
    $$f(x) = x^{1/2} + x^{-1/2}.$$
    $$f'(x) = \frac{1}{2} x^{-1/2} + \left(-\frac{1}{2}\right) x^{-3/2} = \frac{1}{2} x^{-1/2} - \frac{1}{2} x^{-3/2}.$$
    写成根式形式：
    $$f'(x) = \frac{1}{2\sqrt{x}} - \frac{1}{2x\sqrt{x}} = \frac{\sqrt{x}}{2x} - \frac{1}{2x\sqrt{x}} = \frac{x-1}{2x\sqrt{x}}.$$
    \textbf{答案：} $f'(x) = \dfrac{1}{2\sqrt{x}} - \dfrac{1}{2x\sqrt{x}}$ 或 $\dfrac{x-1}{2x\sqrt{x}}$。

    \item 已知 $y = a^x$（$a>0$），用导数的定义证明 $(a^x)' = a^x \ln a$。
    
    \textbf{解：} 回到导数的定义，利用指数运算性质和重要极限 $\displaystyle\lim_{\Delta x \to 0}\frac{a^{\Delta x} - 1}{\Delta x} = \ln a$。
    \begin{align*}
        f'(x) &= \lim_{\Delta x \to 0} \frac{a^{x+\Delta x} - a^x}{\Delta x} \\
        &= \lim_{\Delta x \to 0} \frac{a^x (a^{\Delta x} - 1)}{\Delta x} \\
        &= a^x \cdot \lim_{\Delta x \to 0} \frac{a^{\Delta x} - 1}{\Delta x}.
    \end{align*}
    现在证明 $\displaystyle\lim_{\Delta x \to 0} \frac{a^{\Delta x} - 1}{\Delta x} = \ln a$。
    
    令 $t = a^{\Delta x} - 1$，则当 $\Delta x \to 0$ 时 $t \to 0$，且 $a^{\Delta x} = 1 + t$。两边取自然对数得 $\Delta x \ln a = \ln(1+t)$，所以 $\Delta x = \dfrac{\ln(1+t)}{\ln a}$。于是
    $$\frac{a^{\Delta x} - 1}{\Delta x} = \frac{t}{\frac{\ln(1+t)}{\ln a}} = \ln a \cdot \frac{t}{\ln(1+t)} = \ln a \cdot \frac{1}{\frac{\ln(1+t)}{t}}.$$
    当 $t \to 0$ 时，由重要极限 $\displaystyle\lim_{t \to 0}\frac{\ln(1+t)}{t} = 1$，因此
    $$\lim_{\Delta x \to 0} \frac{a^{\Delta x} - 1}{\Delta x} = \ln a \cdot \frac{1}{1} = \ln a.$$
    代入原式得：
    $$f'(x) = a^x \cdot \ln a.$$
    \begin{quote}
        \textbf{结论：} $(a^x)' = a^x \ln a$，特别地，当 $a=e$ 时，$(e^x)' = e^x$。
    \end{quote}
\end{enumerate}

\section*{写在最后}
亲爱的同学，如果你已经顺利完成了以上练习，恭喜你！你已经熟练掌握了导数的基本计算法则。接下来，你可以将目光转向导数的应用：切线方程、函数的单调性与极值、最值问题等。这些内容才是高考与各类考试的重点。不妨找一本教辅资料，从“导数的定义”跨越到“导数的应用”，开启新的学习旅程吧！

祝你学习进步，数学之路越走越宽广！

\end{document}